%% I. Giới thiệu: bối cảnh, vấn đề, khoảng trống, câu hỏi nghiên cứu, đóng góp, bố cục. Được access_vi.tex \input.
\section{Giới thiệu}
\label{sec:intro}
Nghiên cứu về thiền đã chuyển từ khác biệt trung bình giữa người thiền và nhóm đối chứng sang mô hình hoá \emph{động lực} của một phiên: sự luân phiên giữa chú ý tập trung, lang thang tâm trí và siêu nhận thức \cite{Hasenkamp2012,Lutz2008,Christoff2016}, sự điều tiết mạng chế độ mặc định (DMN) \cite{Raichle2001,Brewer2011,Garrison2015}, và các cách giải thích tính toán theo xử lý dự đoán \cite{Laukkonen2021,SandvedSmith2021}. Khung lý thuyết Bộ não như một giao diện truy vấn (Brain as a Query Interface, BQI) của Vũ \cite{BQI} viết động lực đó thành phương trình vi phân nối hoạt động DMN, thông tin tích hợp $\Phi$ và chất lượng ý thức (QoC) do người tập báo cáo, fit các tốc độ và hệ số khuếch đại từ dữ liệu phân tích gộp, và đề xuất điều khiển chúng bằng một giao diện não--máy vòng kín \cite[Mục~7.2, 9.3 và 10.1]{BQI}; đây là nền tảng lý thuyết của nghiên cứu này (Mục~\ref{sec:theory}). Các mô hình kiểu này hứa hẹn những đại lượng ước lượng được từ dữ liệu, như tốc độ DMN bị ức chế và trọng số của mỗi trạng thái ẩn trong báo cáo, và được đưa ra làm nền cho neurofeedback dựa trên mô hình \cite{Sitaram2017}.

Liệu các đại lượng đó có ước lượng được từ loại dữ liệu mà nghiên cứu thiền thực sự thu, EEG liên tục và vài lần hỏi tự chấm mỗi phiên, là một câu hỏi riêng. Trong bài đồng hành \cite{companion}, chúng tôi thay hệ phương trình của BQI bằng một mô hình không gian trạng thái tuyến tính--Gauss (M1) và dùng thông tin Fisher kỳ vọng chính xác của nó để tính cần bao nhiêu người, phiên dài bao nhiêu, hỏi dày bao nhiêu và đặc trưng EEG sạch đến đâu để ước lượng mỗi tham số đến một độ chính xác cho trước. Phân tích thiết kế đó mới chỉ được kiểm chứng trên dữ liệu mô phỏng. Chưa nghiên cứu nào kiểm phân tích thiết kế như vậy trên bản ghi thật, so các ước lượng với giá trị mà BQI rút ra từ tài liệu, hay hỏi các tham số ước lượng được có chuyển được từ phòng thí nghiệm này sang phòng thí nghiệm khác không (Mục~\ref{sec:gap}).

Trả lời sạch các câu hỏi này không dễ. Các đặc trưng EEG đứng thay cho DMN và cho $\Phi$ là đại diện chưa có ánh xạ hiệu chuẩn sang biến trạng thái; một phần ba bản ghi có thể bị thiếu; mô hình có thể sai đặc tả theo cách mà một cận không lường trước được; và một phân tích có nhiều lựa chọn tự do (đặc trưng, quy tắc nhiễu giả, mô hình rỗng, tiêu chí loại) có thể hậu thuẫn gần như mọi kết luận sau khi đã xem dữ liệu. Vì vậy chúng tôi đăng ký trước toàn bộ kế hoạch phân tích cho mỗi bộ trong hai bộ dữ liệu mở, khoá bằng mã băm SHA-256 trước khi xem EEG \cite{Nosek2018}, tổ chức pipeline thành một chuỗi cổng kiểm tra tự động từ chối chạy tiếp khi điều kiện tiên quyết không đạt, và báo cáo mọi phân tích đã khoá dù có ủng hộ dự đoán hay không.

\subsection{Câu hỏi nghiên cứu}
\begin{enumerate}
  \item[RQ1] \emph{Độ chính xác.} Trên một thiết kế lấy mẫu trải nghiệm thật, với lịch hỏi không đều, epoch thiếu và loại nhiễu giả, sai số tương đối của tham số khó nhất, hằng số tốc độ DMN $\kinh$, có gần cận Cramér--Rao (CRB) tính từ lịch hỏi thật trước khi xem EEG không?
  \item[RQ2] \emph{Dự đoán.} Một mô hình xác định được về nguyên tắc, fit lên EEG liên tục và các báo cáo thưa, có dự đoán báo cáo ngoài mẫu tốt hơn các mô hình bỏ qua động lực không?
  \item[RQ3] \emph{Khớp với BQI.} Các hằng số tốc độ, hệ số ghép và hệ số khuếch đại ước lượng được có khớp với các giá trị mà BQI fit từ dữ liệu phân tích gộp không?
  \item[RQ4] \emph{Lặp lại.} Các tham số ước lượng được có cùng nghĩa giữa các bản ghi của các phòng thí nghiệm khác nhau, với bộ khuếch đại, montage, lịch hỏi và thang điểm khác nhau không?
\end{enumerate}

\subsection{Đóng góp}
\begin{enumerate}
  \item Phép thử đầu tiên, theo hiểu biết của chúng tôi, của một phân tích thiết kế bằng thông tin Fisher cho mô hình động lực của thiền trên EEG lấy mẫu trải nghiệm thật, trên hai bộ dữ liệu mở độc lập, ds001787 \cite{Brandmeyer2018,Delorme2024ds} và OSF 3uszv \cite{RodriguezLarios2021}, theo hai bản đăng ký trước đã khoá và tám cổng kiểm tra.
  \item Một cách báo cáo ba con số sai số, CRB tính trước khi khoá, sai số chuẩn Fisher tại ước lượng và sai số bootstrap tham số, cùng với sai số theo thông tin quan sát và so sánh với dữ liệu mô phỏng từ fit. Trên cả hai bộ, ba con số đầu khớp nhau với hằng số tốc độ DMN, và con số thứ tư để lộ sai đặc tả của kênh công suất alpha \cite{White1982}.
  \item Một kiểm định ngoài mẫu đã đăng ký trước của mô hình so với ba mô hình rỗng, cho thấy không đại diện EEG nào mang thông tin về báo cáo ngoài phụ thuộc chuỗi của chính báo cáo, trên cả hai bộ.
  \item Một phép so sánh các ước lượng với giá trị tham số của BQI, cùng các kiểm tra lặp lại đã đăng ký trước giữa hai bộ, cho thấy các hằng số tốc độ và khác biệt nhóm giữa người thiền lâu năm và người mới không chuyển được từ bộ này sang bộ kia, kèm một mô phỏng giải thích vì sao: fit Ornstein--Uhlenbeck trả về thang thời gian trội của đặc trưng được đưa vào.
  \item Một pipeline tái tạo được hoàn toàn (mỗi bước một lệnh, có checkpoint, 114 unit test) cùng cả hai bản đăng ký trước, khoá và hồ sơ các cổng.
\end{enumerate}

Phần còn lại của bài theo bố cục thông thường của một bài nghiên cứu. Mục~\ref{sec:litreview} tổng quan tài liệu về mô hình động lực của thiền, khung BQI, các chỉ dấu EEG của hoạt động DMN và của tích hợp, khả năng xác định và thiết kế thí nghiệm cho mô hình động lực, và đăng ký trước, rồi nêu khoảng trống nghiên cứu. Mục~\ref{sec:theory} trình bày khung lý thuyết: mô hình BQI với các phương trình của nó, dẫn xuất mô hình ước lượng được M1 và các dự đoán của BQI cho các đại lượng được ước lượng. Mục~\ref{sec:methods} trình bày phương pháp: thiết kế nghiên cứu, mô hình, giả thuyết, dữ liệu, đăng ký trước, tiền xử lý và phân tích thống kê. Mục~\ref{sec:results} báo cáo kết quả, Mục~\ref{sec:discussion} thảo luận chúng cùng các hạn chế của nghiên cứu, và Mục~\ref{sec:conclusion} kết luận và nêu hướng tiếp theo.
